%Note for this test case:
%Simple water flow example using ANUGA: Water flowing down a channel.
%It was called "steep_slope" in an old validation test.


\section{Shallow flow down a mild slope}
This case simulates very shallow flow running down a mild slope topography. It represents an idealisation of the rainfall-runoff problem, which will often involve very shallow flows down such a topography. This case has an analytical solution, and in particular, we consider the steady-uniform solution with the values of bed slope and friction slope are the same.   

Suppose that we are given a one dimensional domain. The steady state conditions with a contant water depth everywhere make the shallow water equations to the single identity
\begin{equation}
z_x = - S_f, 
\end{equation}
where $z_x$ is the bed slope, and $S_f$ is the symbol for the force of bottom friction.  We take Manning's friction
\begin{equation}
S_f = n^2 \frac{q|q|}{h^{10/3}}
\end{equation}
where $n$ is the Manning's coefficient and $q$ is the discharge $uh$. 
If $q$, $n$, and $z_x$ are given, then the analytical solution for $u$ and $h$ is
\begin{equation}
u(x)= \left[- n^{-2} q^{4/3} z_x\right]^{3/10},
\end{equation}
\begin{equation}
h(x)= \frac{q}{u}\,.
\end{equation}

\subsection{Stability of the uniform flow: roll waves}
The steady uniform flow is not always stable. When the Froude number
$F = u/\sqrt{gh}$ exceeds a threshold -- $3/2$ for Manning friction, $2$ for
Ch\'ezy friction -- small disturbances grow as they travel down the slope and
form \emph{roll waves}: a train of bores separated by thin, fast
flow~\cite{Dressler1949}. The growth is physical, not numerical, so a scheme
that holds the uniform state in this regime is too dissipative rather than
more accurate. Reducing the CFL number, or using a more diffusive scheme,
delays the onset of roll waves without removing them.

We therefore run two cases with the same discharge $q=0.2$ and bed slope
$z_x=-0.1$:
\begin{itemize}
\item a \emph{stable} case with $n=0.06$, so that $u=1.424$, $h=0.1405$ and
  $F=1.21$, on a $100 \times 100$ domain. The numerical solution should
  reproduce the analytical one;
\item a \emph{roll-wave} case with $n=0.03$, so that $u=2.158$, $h=0.0927$ and
  $F=2.26$, on a channel $200$ long and $20$ wide. The uniform flow is
  unstable here, and the numerical solution should depart from it
  downstream.
\end{itemize}
Both use a $2$ m grid of cross-triangulated squares. The topography is
\begin{equation}
z(x, y)= -0.1 x\,.
\end{equation}
The initial condition is $u=v=0$ and
\begin{equation}
w(x,y,0)= -0.1 x + 0.01\,.
\end{equation}
The upstream boundary imposes the analytical depth and discharge, the
downstream boundary is transmissive and the sides are reflective walls.
In the roll-wave case the inflow discharge carries a small disturbance,
\begin{equation}
q_{\mathrm{in}}(t) = q \left(1 + 0.01 \sin \frac{2\pi t}{3.3}\right),
\end{equation}
so that every flow algorithm is given the same seed. Without it, a scheme that
settles to an exact discrete steady state has nothing to amplify and shows no
roll waves, however weakly it damps them; the outcome would then depend on
whether the scheme happens to generate its own unsteadiness rather than on
how it treats the instability.

\subsection{Results: stable case}
Figure~\ref{fig:depthdownchan} shows the steady state depth in the downstream direction.
There should be a good agreement with the analytical solution, at least away from the boundaries.

Figure~\ref{fig:depthacrosschan} shows the steady state depth across the slope around the line x = 50m.
There should be a good agreement with the analytical solution, at least away from the boundaries.

Figures~\ref{fig:xvelscrosschan} and~\ref{fig:yvelscroschan} show the steady state $x$- and $y$-velocities, along a slice in the cross slope direction (near $x=50$). The $x$-velocities should agree well with the analytical solution, and the $y$-velocities should be zero.

\begin{figure}
\begin{center}
\includegraphics[width=0.8\textwidth]{depth_x.png}
\caption{Stable case: depth in the downstream direction}
\label{fig:depthdownchan}
\end{center}
\end{figure}

\begin{figure}
\begin{center}
\includegraphics[width=0.8\textwidth]{depth_y.png}
\caption{Stable case: depth across the slope around $x$ = 50m}
\label{fig:depthacrosschan}
\end{center}
\end{figure}

\begin{figure}
\begin{center}
\includegraphics[width=0.8\textwidth]{x_velocity.png}
\caption{Stable case: $x$-velocity along the cross-section $x=50$ (i.e. a cross-section with constant bed elevation)}
\label{fig:xvelscrosschan}
\end{center}
\end{figure}

\begin{figure}
\begin{center}
\includegraphics[width=0.8\textwidth]{y_velocity.png}
\caption{Stable case: $y$-velocity along the cross-section $x=50$ (i.e. a cross-section with constant bed elevation)}
\label{fig:yvelscroschan}
\end{center}
\end{figure}

\subsection{Results: roll-wave case}
There is no steady solution to compare with here. What we look for is the
behaviour of the instability: the flow should stay close to uniform near the
inflow, where the 1\% disturbance is small; the disturbance should then grow
exponentially with distance down the slope; and it should saturate into roll
waves of finite amplitude.

Figure~\ref{fig:rollwavedepth} shows the depth along the centreline at the
end of the run, and its range over the last quarter of the run.
Figure~\ref{fig:rollwaveamp} shows half that range relative to the uniform
depth: a straight line on the logarithmic scale is exponential growth, and
the plateau is saturation. Figure~\ref{fig:rollwavext} shows the centreline
depth in space and time over the last 30 s of the run, which is stored every
0.5 s: each diagonal stripe is one roll wave travelling down the slope.

The flow is supercritical, so both characteristics point downstream and the
transmissive outflow boundary cannot influence where the waves start. Their
onset is set by the size of the inflow disturbance and the rate at which it
grows.

\begin{figure}
\begin{center}
\includegraphics[width=0.8\textwidth]{rollwave_depth_x.png}
\caption{Roll-wave case: depth along the centreline, at the final time and its range over the last quarter of the run}
\label{fig:rollwavedepth}
\end{center}
\end{figure}

\begin{figure}
\begin{center}
\includegraphics[width=0.8\textwidth]{rollwave_amplitude.png}
\caption{Roll-wave case: amplitude of the depth fluctuation down the slope, relative to the uniform depth}
\label{fig:rollwaveamp}
\end{center}
\end{figure}

\begin{figure}
\begin{center}
\includegraphics[width=0.8\textwidth]{rollwave_xt.png}
\caption{Roll-wave case: centreline depth relative to the uniform depth, in space and time, over the last 30 s}
\label{fig:rollwavext}
\end{center}
\end{figure}


\FloatBarrier
\subsection{Comparison of flow algorithms}
The roll-wave case is run with each flow algorithm (DE0, DE1, DE2 and
DE\_ader2, at their default settings) by \texttt{compare\_solvers.py}. All of
them hold the stable case, so this is where they differ. The physical growth
is the same for every scheme, so differences in where the disturbance reaches
a given amplitude measure the scheme's numerical damping.
Figure~\ref{fig:rollwavesolvers} compares the growth of the wave amplitude,
and Figure~\ref{fig:rollwavesolversxt} the centreline depth in space and
time over the last 30 s. The second-order schemes should grow the disturbance at nearly the same
rate. DE0, with first-order timestepping and limiter $\beta = 0.5$, is
expected to damp it much more strongly.

\begin{figure}[htb]
\begin{center}
\includegraphics[width=0.8\textwidth]{rollwave_solvers.png}
\caption{Roll-wave case: growth of the wave amplitude down the slope for each flow algorithm}
\label{fig:rollwavesolvers}
\end{center}
\end{figure}

\begin{figure}[htb]
\begin{center}
\includegraphics[width=\textwidth]{rollwave_solvers_xt.png}
\caption{Roll-wave case: centreline depth relative to the uniform depth, in space and time over the last 30 s, for each flow algorithm}
\label{fig:rollwavesolversxt}
\end{center}
\end{figure}


\endinput
